\section{Computations}\label{app:comp}

This appendix describes the finite computations. Every result marked as computer-assisted in the text is the output of a terminating exact search over a finite set; where random sampling is used, this is said, and such runs serve only as checks. The directory \texttt{code/} contains the programs, with instructions for running them and their expected outputs, and the logs of the runs; the programs, or their directories, are named below next to each computation. Solution tables and other data files of more than $5$\,MB are not included; the programs that produce them are. Section~\ref{app:census} describes the census of Lemma~\ref{cl:lem:95cens}, Section~\ref{app:stated} the other computations whose results are stated in the text without a proof by hand, and Section~\ref{app:confirm} computational confirmations of results that are proved by hand.

\subsection{The census for rule 95}\label{app:census}

Let $g$ be a block map of radius $R$ with local rule $\gamma$, and $F=F_{95}$. Both $g\circ F$ and $F\circ g$ are block maps of radius $R+1$, so they are equal if and only if their local rules agree on every window $w\in\A^{2R+3}$. The unknowns are the $2\cdot4^{2R+1}$ bits of the table of $\gamma$. For a window $w$ let $e_1,e_2,e_3$ be its three subwords of length $2R+1$, and $e_L$ the word of length $2R+1$ obtained by applying the local rule of $F$ to $w$; writing $X_e,Y_e$ for the two bits of the table entry $e$, the window contributes the two constraints
\[
X_{e_L}=Y_{e_2},\qquad Y_{e_L}\oplus X_{e_2}=f(Y_{e_1},Y_{e_2},Y_{e_3}),
\]
which are encoded as clauses by their truth tables. All solutions are enumerated by repeated solving with blocking clauses, until the instance is unsatisfiable, using PySAT \cite{PySAT} (\texttt{code/census/centralizer.py} for $R\le3$, \texttt{enum\_r4.py} for $R=4$). Each enumeration is carried out twice, with the solvers CaDiCaL~1.5.3 \cite{CaDiCaL} and MiniSat~2.2 \cite{MiniSat}, and for every $R\le4$ the two solvers return the same set of tables. Every table returned by the solver is checked again, independently of the clauses, by evaluating $g\circ F$ and $F\circ g$ on all $4^{2R+3}$ windows. It is then compared with the tables of the maps $\sh^aF^b$ with $|a|+|b|\le R$, which are built by composing local rules, and with the constant tables; the program checks that every $\sh^aF^b$ with $|a|+|b|\le R$ occurs among the solutions. The instances have $128$, $2048$, $32768$ and $524288$ variables and $17542$, $292998$, $4714182$ and $75486342$ clauses for $R=1,2,3,4$, and $6$, $14$, $46$ and $262$ solutions: $5$, $13$, $25$ and $41$ maps $\sh^aF^b$, one constant table, and $0$, $0$, $20$ and $220$ others (the four logs \texttt{census\_95\_R123\_*.txt} and \texttt{census\_95\_R4\_*.txt} in \texttt{code/census/logs}, one for each solver). That an enumeration is complete rests on the final answer of the solver that the instance with all blocking clauses is unsatisfiable. These answers are not certified by proof logs; the evidence for them is the agreement of the two solvers, and, for $R\le4$, the agreement with the search of Section~\ref{app:95} over tables of sublattice type, which has a different encoding. The same holds for the other searches with a SAT solver in this appendix. The solutions themselves do not depend on the solver, since each is verified directly. For $R\le3$ the program also decides the injectivity of each solution on the graph of pairs of words: the injective ones are the maps $\sh^aF^b$. This test is not used in a proof here, since the further maps are identified in Proposition~\ref{cl:prop:95census}(a). The same programs are applied to all $88$ classes in \cite{PaperD}.

\subsection{Computer-assisted statements}\label{app:stated}

The following results are stated in the text as computations. For each we give the method in brief.

\subsubsection{Rule $155$: window closures and front strips}\label{app:closure} (Remarks~\ref{cl:rem:155} and~\ref{cl:rem:155G1}, Lemma~\ref{cl:lem:155strip}.) The \emph{window closure} of a rule is a monotone closure on the set of windows $w\in\A^{2R+1}$ at which a block map $h$ commuting with $F$ is known to act as the identity. It starts from the windows of configurations that $h$ fixes by hypothesis, and it uses the locally permutive readings of the spacetime relation for $\Delta=Y_{hc}\oplus Y_c$: a window is added if, for some extension of it by two cells on each side and one of the five instances of the relation through the centre, the relation has no solution with $\Delta=1$ at the centre and $\Delta=0$ at those of its other four points whose windows are already in the set, the remaining values being arbitrary; this includes the instances in which the rule does not depend on an argument in the given context. Words of $\Phi_f$ are admitted to the base if they extend by six letters on each side. For $155$ the base consists of the windows of $\Phi^+\cup\Phi^-$, and the admitted words are those that extend to bi-infinite points; the closure is complete for $R\le6$ (\texttt{code/closure/dc.c} for $R\le5$; \texttt{code/rule155/mkbase.py} and \texttt{lc.c} for $R=6$, with the log \texttt{lc155\_R6.log}). The weaker closure of Remark~\ref{cl:rem:155G1} uses only the reading that determines the value at the point $\mathsf S$ (\texttt{sonly.py} for $R=2,3$; \texttt{mkseed.py} and \texttt{sc.c} for $R=4$). With the base $\Phi^+\cup\Phi^-$ alone it leaves $968$, $16140$ and $261168$ windows for $R=2,3,4$. It is complete for $2\le R\le4$ when the base also contains the windows of the periodic front strips of both kinds of front, centred at the offsets $-R\le s\le R$ from the edge ($-4\le s\le5$ for $R=4$), and of their time reversals; for $R=2$ and offsets $-2\le s\le1$ it leaves $875$ windows. The periodic continuations of Lemma~\ref{cl:lem:155strip}(i) are enumerated line by line from \eqref{cl:eq:strip}, by classifying each line as slaved, free or doubling (\texttt{strip.py}; \texttt{s4.py} for the $00$-front at width $18$).

\subsubsection{Rule $91$}\label{app:91} (Remark~\ref{cl:rem:91} and Lemma~\ref{cl:lem:91Q1}.) The objects are the points of $\Phi^-$ on the background $11$ whose core, a word of $\Phi_{91}$ that begins and ends with $0$, has length at most $12$; there are $778$. For each of them the scattering of a bare defect $\Delta^+_a$ off it is computed at both relative parities, and the object counts as crossed unchanged if $\mathcal S(\Delta^+_a,X)=(X,\Delta^+_a)$ at both. Two programs do this (\texttt{code/rule91/t91b.py} and \texttt{t91c.py}), one with the recurrence \eqref{cl:eq:goursat} and one by direct simulation of the cells in $(t,i)$, in which the configuration after the crossing is compared with the free superposition $\pi(q)\oplus\xi(p)\oplus1$; both find that the objects crossed unchanged are exactly the $41$ whose core is $3$-sparse. The two example crossings with the object $1^\infty001^\infty$ at the end of the first paragraph of the remark are computed with the recurrence \eqref{cl:eq:goursat} (\texttt{s91.py}). For the conserved quantities, the additive conserved quantities of $F_f$ with a density $h$ on $r$ adjacent cells are the solutions of the linear system $h(F(w))-h(w_1\cdots w_r)=g(w_1\cdots w_{r+1})-g(w_0\cdots w_r)$, $w\in\A^{r+2}$, in the unknowns $h$ and $g$, where $F(w)$ is the image of the $r$ inner cells; the system is solved exactly over $\mathbb Q$ (\texttt{cq.py}, \texttt{cq\_check.py}). That the densities of $Q_1$ and $Q_2$ solve it for $r=2$ is an exact rank computation, for $Q_2$ with $f\in\{91,94,95\}$ and for $Q_1$ with $f=91$; for $94$ and $95$ the density of $Q_1$ does not solve it. Both quantities are also evaluated on $c$ and $F(c)$ for all configurations $c$ of spatial period at most $7$ (\texttt{q91.py}). These quantities are known \cite{Takesue1995,Takesue1997}; the computation is a check of the formulas in the present notation. The identity in the proof of Lemma~\ref{cl:lem:91Q1} is checked on all cases as well: the regrouping of the density, the cases ($\alpha$) and ($\beta$), and the invariance of each term (\texttt{q1\_identity.py}).

\subsubsection{The statement \textup{(Sh)}}\label{app:sh} (Remark~\ref{cl:rem:IR}; program in \texttt{code/sh}.) All scatterings are computed by direct simulation in $(t,i)$. For each object with a core of the stated length and each pair of tails, the probes with cores up to the stated length are searched for one that shortens the object (\texttt{sh.c}; $452060$ and $793310$ probes with cores of length at most $24$ on the two backgrounds of $94$, and $167914$ and $203804$ with cores of length at most $20$ on the backgrounds $01$ and $11$ of $91$, each core being counted at both of its placements modulo $2$; the background $10$ of $91$ is the image of $01$ under the shift and has the same number of probes; log \texttt{logs/sh\_all.log}). For an object with two different tails only those probes are used that send the two bare kinks with these tails to bare kinks, so the search establishes more than (Sh) in the stated range.

\subsubsection{Rule $95$: the smallest radius and the commuting maps of small radius}\label{app:95} (Lemma~\ref{cl:lem:95enum}, Propositions~\ref{cl:prop:95small} and~\ref{cl:prop:95census}; the programs, the tables and the logs are in \texttt{code/rule95}.)
\emph{Particle-conserving block maps.} For $w\le8$ the functions $\lambda$ of the criterion after Lemma~\ref{cl:lem:95enum} are enumerated by a depth-first search over the vertices of $\mathsf B_w$ in breadth-first order from $0^{w-1}$; a value $\lambda(v)$ has at most two candidates once a neighbour of $v$ is assigned, and every edge and every pair of consecutive edges is checked as soon as its vertices are assigned. The program \texttt{enum\_golden.py} gives the numbers of Table~\ref{cl:tab:95enum} (log \texttt{logs/enum\_golden\_8.log}). For $w\le7$ a second computation (\texttt{enum\_golden\_linear.py}) imposes conservation as the linear equations $\sum_\epsilon m_\epsilon(z)\hat\varphi(\epsilon)=|z|$ over all cyclic words $z$ of $\mathsf G$ of length at most the number of vertices, where the sum is over the edges $\epsilon$ of $\mathsf B_w$ and $m_\epsilon(z)$ counts the occurrences of $\epsilon$ in $z$ (these include all simple cycles of $\mathsf B_w$), reduces them over $\mathbb Q$, and tries all $0/1$ values of the free coordinates ($2^{20}$ for $w=7$); it gives the same sets, and the numbers for $w=7,8$ are obtained again by the SAT search below. Injectivity is decided on the graph whose vertices are the pairs of vertices of $\mathsf B_w$ and whose edges are the pairs of edges with equal values of $\hat\varphi$: the map is injective if and only if no vertex off the diagonal lies on a bi-infinite path, that is, survives the repeated removal of vertices without a successor or without a predecessor. For each of the $469$ maps with $w=7$ and the $13506$ maps with $w=8$ that are not injective, two distinct points of $\mathsf G$ of the same period, at most $10$, with the same image are found as well (\texttt{inj\_witness.py}).

\emph{The automorphism of radius $7$.} The map $\widehat{H_{\varphi,\psi}}\sh^n$ is evaluated on a finite perturbation $c$ of $b_{11}$ from the properties that characterise it: $F^{-t}$ is applied until the configuration consists of a point of $\Phi^+$ on the left and a point of $\Phi^-$ on the right, separated by at least $64$ cells $(1,1)$; there $\varphi$ is applied to the two classes of zeros of the left part and $\psi$ to those of the right part, as in \eqref{cl:eq:95restr}; then $F^t$ and $\sh^n$ are applied. With $(\varphi,\psi)=(\sh,\mathrm{id})$, $(\mathrm{id},\sh)$ and $(\sh^{-1},\sh)$ this reproduces $\sh^{-1}F$, $\sh F$ and $\sh^2$ on $300$ pseudo-random finite perturbations each (the self-test of \texttt{core.py}). The tables $T_x,T_y$ of Proposition~\ref{cl:prop:95small}(b) are filled by evaluating $h_{\ast}$ on configurations in which the two classes of bits run through a binary de Bruijn sequence of order $15$, so that every entry is hit, and every sample is required to be consistent with the entries already filled (\texttt{gen7.py}). The table of $h_\ast$ found by the SAT search of the paragraph on tables of sublattice type below is identical. The four facts of the proof are then checked on the tables alone: the commutation on $2\cdot2^{17}$ windows and the involution on $2\cdot2^{29}$ windows by a program that evaluates $F$ and the table map cell by cell, with the bits of the other class drawn at random, and in addition on $2^{26}$ windows with all bits random (\texttt{verify7.c}; logs \texttt{verify7.log} and \texttt{verify7\_negative\_control.log} in \texttt{code/rule95/logs}); the restrictions on the $3025$ words of length $16$ of $\Phi_{95}$; and the dependence by comparing every entry with the $15$ entries at Hamming distance $1$ (\texttt{props7.py}). With one table entry complemented, the commutation check reports $12$ mismatches and the involution check $49152$. The commutation and the involution are also confirmed on all configurations of spatial period at most $9$; $\sh h_{\ast}$ moves $310$ of those of period $5$ and $7168$ of those of period $7$, and for each of these $h_{\ast}(c)$ is not in the orbit of $c$ under $\langle\sh,F\rangle$ (\texttt{props7.py}; \texttt{standalone7.py} for the configuration of period $5$ in the proof).

\emph{Tables of sublattice type.} For a radius $R$ and each of the two ways in which the two bits of the image can be assigned to the two classes, the unknowns are the $2\cdot2^{2R+1}$ entries of two tables, and $hF=Fh$ at the cell $0$ consists of $2^{2R+3}$ equalities between two entries and $2^{2R+3}$ relations $a=b\oplus1\oplus cd$ between four entries. For $R\le5$ all solutions are enumerated with a SAT solver and blocking clauses (\texttt{sat\_sublattice.py}; CaDiCaL~1.5.3, and a second time with MiniSat~2.2, with the same sets of solutions); the numbers of solutions are $2+5$, $10+5$, $10+37$, $226+37$ and $226+2917$ for $R=1,\dots,5$, the constant map being counted in both types, which for $R\le4$ are the totals $6$, $14$, $46$, $262$ of the census of Lemma~\ref{cl:lem:95cens}, made without this restriction. For $R\le7$ the restrictions to $\Phi^+$ that occur among the solutions are enumerated by blocking on the entries at windows of $\Phi^+$ (at most $110$ variables), likewise for $\Phi^-$, the injective ones are selected by the test above, and for every pair of an injective restriction to $\Phi^+$ and one to $\Phi^-$ all solutions with these restrictions are enumerated. An automorphism of sublattice type is among them, since its restrictions to $\Phi^\pm$ are injective. The numbers of restrictions to $\Phi^+$ (and to $\Phi^-$) that occur are $\mathsf m_w+1$, with $\mathsf m_w$ from Table~\ref{cl:tab:95enum} and $w$ as in Lemma~\ref{cl:lem:95window}(ii), the additional one belonging to the constant map: $477$ for the type without exchange of the classes at $R=6$ and $R=7$, and $55$ and $13517$ for the type with exchange; the injective ones number $7$, $7$, $6$ and $10$, and the class maps of all of them satisfy the criterion after Lemma~\ref{cl:lem:95enum}. This reproduces the last two columns of Table~\ref{cl:tab:95enum} by a different method. For $R\le6$ the solutions with injective restrictions are the $2R^2+2R+1$ maps $\sh^aF^b$; for $R=7$ they are the $113$ maps $\sh^aF^b$ and $36$ others, all of the type with exchange. The tables of the $36$ maps of Proposition~\ref{cl:prop:95small}(c) are also generated from the definition, as for $h_{\ast}$ (\texttt{gen36.py}), and verified as in the proof of that proposition (\texttt{verify36.py}: commutation on all $2\cdot2^{17}$ windows, restrictions on all words, dependence on the cells $i\pm7$; \texttt{sat\_post.py} writes the tables found by the search); the two sets of $36$ pairs of tables are identical. The instance for $R=7$ has $65536$ variables and about $1.2\cdot10^6$ clauses. The searches for $R\le7$ are carried out with both solvers, with the same sets of solutions; the logs are in \texttt{code/rule95/logs}, those of the search for $R=7$ being \texttt{sat\_sublattice\_aut\_R7.log} and \texttt{sat\_sublattice\_aut\_R7\_minisat22.log}.

\emph{Census tables.} For each of the $46$ tables of the census of radius $3$ and the $262$ of radius $4$ (Lemma~\ref{cl:lem:95cens}), the program \texttt{census\_cmp.py} rechecks the commutation on all $4^{2R+3}$ windows, verifies on all $4^{2R+1}$ entries that the table is of sublattice type, extracts the four class maps of its restrictions to $\Phi^+$ and $\Phi^-$ from the entries at the windows of $\Phi^\pm$, checks that they have the form \eqref{cl:eq:95restr} of $\widehat{H_{\varphi,\psi}}\sh^n$ with $n\in\{0,1\}$ and that $\varphi$ and $\psi$ satisfy the criterion after Lemma~\ref{cl:lem:95enum}, and decides their injectivity. As a second check, the table of $\widehat{H_{\varphi,\psi}}\sh^n$ is generated by the evaluation described above and compared with the census table in all entries. All $45$, resp.\ $261$, non-constant tables pass both checks; $25$, resp.\ $41$, have $\varphi$ and $\psi$ injective, and these are the maps $\sh^aF^b$. The same two checks are applied to the $3141$ non-constant solutions of sublattice type of radius $5$ (\texttt{verify\_sub.py}), and all pass. The census tables of radius $3$ are in \texttt{code/rule95}; the $262$ tables of radius $4$, a file of $69$\,MB, are not included and are produced by \texttt{code/census/enum\_r4.py}.

\subsubsection{Rule $175$}\label{app:175} (Remark~\ref{cl:rem:175}.) The program \texttt{sat\_sublattice.py} takes the rule as a parameter: for an elementary rule $f$ that does not read its middle argument, the relations between four entries are $a=b\oplus f(c,\cdot,d)$. For $f=175$ and each $R\le7$ all solutions of both types are enumerated: they are the $2R^2+2R+1$ maps $\sh^aF^b$, with no constant map (no constant configuration is fixed by $F_{175}$) and no other solution, and every solution is checked again to commute with $F_{175}$ on random periodic configurations. A second program with its own encoder (\texttt{sublattice\_any.py}) gives the same numbers of clauses and of solutions for every $R\le7$. Both programs are run with CaDiCaL~1.5.3 and with MiniSat~2.2, with the same results. As a control, for $f=250$ and $R=3$ the first program finds the $25$ maps $\sh^aF^b$, the constant map, counted in both types, and $20$ others, as for $f=95$ (log \texttt{logs/sat\_sublattice\_rule250\_R3.log}). The instances for $R=7$ have $65536$ variables and about $1.2\cdot10^6$ clauses for each type. The logs of both programs are in \texttt{code/rule95/logs}.

\subsection{Confirmations of results proved by hand}\label{app:confirm}

The following computations check results that are proved by hand. None of them is used in a proof.

\emph{Lemma~L} (Theorem~\ref{cl:thm:L}; programs in \texttt{code/lemmaL} and \texttt{code/rule155}). Lemma~\ref{cl:lem:climb}(i)--(iii) is checked over all words of length at most $22$ for $91$ and $94$ (\texttt{climb.c}); the formula of Lemma~\ref{cl:lem:dalembert} at the $72$ lattice points of a block of $12\times12$ in each of $200$ pseudo-random diagrams per rule, $14{,}400$ points (\texttt{dalembert.py}); and Proposition~\ref{cl:prop:sep} for $91$ and $94$, for all $4^7$ fillings of $7$ adjacent cells in every valid background with $R=4$, all of which separate by $t=18$ (\texttt{separ.py}). Working with \eqref{cl:eq:goursat} directly, the two-step gains of the climbing lemmas are $4$, $3$, $4$ for $91$, $94$, $95$ (\texttt{climb\_gains.c}). The window closure of Section~\ref{app:stated} with base $\Phi^+\cup\Phi^-$ is complete for $R\le4$ for $91$ and $94$ (\texttt{code/closure/dc.c}). For $155$, Proposition~\ref{cl:prop:155} is checked on $3000$ pseudo-random configurations (\texttt{edge155.py}); Proposition~\ref{cl:prop:155W} on all fillings of at most $4$ cells next to the medium with three pseudo-random media each, $1023$ runs of $40$ rows (\texttt{checkW.py}), and on the $99{,}059$ rows of $3000$ pseudo-random runs of at most $60$ steps (\texttt{thmW.py}); and \eqref{cl:eq:strip} and the eventual periodicity of the strips by simulation (\texttt{stripsim.py}).

\emph{Kink rigidity} (Theorems~\ref{cl:thm:KR94} and~\ref{cl:thm:KR95}; programs in \texttt{code/kink}). The following checks use a direct simulation in $(t,i)$ of the configuration that is $P$ on the left and $X$ on the right, independent of \eqref{cl:eq:goursat}. Lemma~\ref{cl:lem:pair} is checked in all eight cases (\texttt{t\_pair.py}). Lemma~\ref{cl:lem:ray} is checked exhaustively for both front values and for the version with $q$ and $p$ exchanged, over all continuations of length at most $10$ of $A$ and $B$, at two values of $a_1$ and two lengths of the front of $B$: $1{,}465{,}860$ cases, with the number and the position of the firings on the ray $a_1-1$ as stated (\texttt{t\_R1.py}; log \texttt{code/kink/logs/t\_R1\_10.log}). Lemma~\ref{cl:lem:dalembert}, in the form used in Step~3, is checked on $800$ random pairs (\texttt{t\_dal.py}), and Lemma~\ref{cl:lem:train} for $2\le g\le8$ and $g\in\{15,22,31\}$ at the four placements $a,b\in\{8,9\}$ (\texttt{t\_train.py}). As an end-to-end check of Steps~2--5 (\texttt{t\_step4.py}), among the $351$ words $w$ of length at most $12$ that are empty or begin with $1$ and end with $0$, and for which $00\,w\,11$ contains no $010$, only the empty word gives a point $W$ (the word $w$ after a front $0^b$, followed by a plateau of $1$s and a bare descending kink at a fixed position) such that the outgoing $\Phi^+$ part of $(K^+_a,W)$ is $K^+_a$ for both $b$ and $b+1$. For $95$, Lemma~\ref{cl:lem:pair95} is checked at both parities (\texttt{t95a.py}), the two parts of Lemma~\ref{cl:lem:rows95} are checked on all pairs of cores of length at most $9$ at two placements and both classes, $6000$ cases each, and among all pairs of cores of length at most $9$ only the pair of bare defects satisfies both relations (T) and ($\Delta$) of the proof of Theorem~\ref{cl:thm:KR95} (\texttt{t95b.py}).

\emph{The translation property for $95$} (Corollary~\ref{cl:cor:95trans}). The property is checked exhaustively for cores of length at most $9$ with \eqref{cl:eq:goursat}, and for cores of length at most $6$ also by direct simulation in $(t,i)$ (\texttt{code/translation95/t95ex.py}: $3200$ placed pairs, resp.\ $200$).

\emph{Rule $95$: blocking words and window graphs} (Propositions~\ref{cl:prop:95noblock} and~\ref{cl:prop:95eigen}; the programs are in \texttt{code/rule95}). For every word $w$ of length at most $6$ and every cell $i_0$ of it, the configurations $c$ and $c'$ of the proof of Proposition~\ref{cl:prop:95noblock}, with the additional particle realised as a bare defect of $\Phi^+$ at distance $4|w|+12$ or $4|w|+13$ to the left, are evolved, and a time with $F^t(c)_{i_0}\ne F^t(c')_{i_0}$ is found in every case, the first such time being at most $47$ (\texttt{block\_test.py}). The graphs $W^{(1,1)}_M$, $W^{(-1,1)}_M$ and $W^{(0,1)}_M$ of $95$R are strongly connected for $M\le7$ (\texttt{window\_graph.py}), in accordance with Proposition~\ref{cl:prop:95eigen}; as a control, $W^{(\pm1,1)}_M$ of $90$R is not strongly connected for $1\le M\le4$, while $W^{(0,1)}_M$ is.

\emph{The scattering closure for $94$} (Lemmas~\ref{cl:lem:pair}, \ref{cl:lem:prepend}, \ref{cl:lem:B}, \ref{cl:lem:exchange} and~\ref{cl:lem:prependZ}, Corollary~\ref{cl:cor:zero}, Theorem~\ref{cl:thm:objects}). Two sets of programs (\texttt{code/rule94}), which share no code, check these results. Sequences are counted up to translation, and the core is that of Section~\ref{cl:sec:scatter}, so that a bare kink has a core of length $0$: there are $153$, $467$, $818$, $1434$ and $2515$ sequences of $\Phi_{94}$ with constant tails and cores of length at most $8$, $10$, $11$, $12$ and $13$, the two backgrounds included.

The first set (\texttt{lib.py}, \texttt{cc.py}, \texttt{verify1.py}, \texttt{verify2.py}, \texttt{verify3.py}, \texttt{closure\_bf.py}) uses the recurrence \eqref{cl:eq:goursat} in two implementations and the direct simulation in $(t,i)$, which agree. It checks the identities $\mathcal S(P_i,[b;U_i])=([b;U_{i+1}],\Omega_i)$ of the two cases of the proof of Lemma~\ref{cl:lem:prepend}, with the bounds on the first kink of $U_{i+1}$, on $3000$ random objects $U_i$ with at most $7$ kinks in each case; Lemma~\ref{cl:lem:B} for $2\le g\le8$ at four placements, $28$ cases; Lemma~\ref{cl:lem:exchange}(i) and (iii) with their firings and outgoing parts, and the outgoing parts of the pairs of (ii), for all $2\le a,d\le11$, and the diagrams of (i) and (iii) on the whole quadrant; and the identities of the two cases of the proof of Lemma~\ref{cl:lem:prependZ} on $3000$ random objects each. It runs the proof of Theorem~\ref{cl:thm:objects} as an algorithm, in which every scattering is computed and compared with the form that the proof asserts, and (C0) is invoked once the prepended kink is at distance $L$ from the rest: with $L=10$ and \eqref{cl:eq:goursat} this handles all $816$ objects that are not backgrounds with cores of length at most $11$ ($11{,}611$ steps of the chains of Lemma~\ref{cl:lem:prepend}), and with $L=8$ and the direct simulation all $465$ objects that are not backgrounds with cores of length at most $10$ ($4467$ steps). Without the lemmas, it computes the closure of the bare points under (C0)--(C3) breadth first on the sequences with cores of length at most $12$ for $L=6$, and at most $8$ for $L=4$, using in (C2) and (C3) the probes with cores of length at most $1$: the closure contains all $1434$, resp.\ $153$, of these sequences.

The second set (\texttt{rc.py}, \texttt{ta.py}, \texttt{tc.py}, \texttt{te.py}) uses only a direct simulation written from the local rule, in which two rows of the configuration that is $P$ on the left and $X$ on the right are evolved on a finite array. It checks Lemma~\ref{cl:lem:pair} and \eqref{cl:eq:cross94}, with the diagram on the quadrant, the firings and the outgoing parts, for all $-4\le a,b\le4$; Lemma~\ref{cl:lem:B}, with its firings, for $2\le g\le12$ in $810$ cases; the identity $X=[b;[r;K^-_{r+2}]]$ used in Corollary~\ref{cl:cor:zero} at $25$ placements; and Lemma~\ref{cl:lem:exchange}(i), (ii) and (iii) in $41$, $41$ and $40$ cases. It checks the identities and the bounds on the first kink of the two cases of the proof of Lemma~\ref{cl:lem:prepend} in $10{,}810$ and $14{,}320$ cases, and those of the two cases of the proof of Lemma~\ref{cl:lem:prependZ} in $11{,}456$ and $2864$ cases; these cover all $2513$ objects that are not backgrounds with cores of length at most $13$, with the prepended kink at distance $2$, $3$, $4$, $5$ and $8$ and the prepended zero at distance $3$, $4$, $5$, $6$ and $9$ from the first kink, and two positions of the probe. The same identities hold for $400$ random objects with cores of lengths $14$ to $40$. The proof of Theorem~\ref{cl:thm:objects} is run as an algorithm in which each chain of Lemma~\ref{cl:lem:prepend} is followed through all its $2L$ scatterings, each scattering $\mathcal S(P_i,[b;U_i])$ is compared with the one computed from $U_i$ alone, and the bound $u_n\ge b+L+2$ and the hypothesis of (C0) are checked position by position: it handles all $818$ sequences with cores of length at most $11$ for $L\in\{1,2,4,6\}$, all $1434$ with cores of length at most $12$ for $L=10$ ($58{,}164$ scatterings), all $2515$ with cores of length at most $13$ for $L=12$ ($122{,}242$ scatterings), all $467$ with cores of length at most $10$ for $L=16$ ($29{,}894$ scatterings), and $1500$, resp.\ $300$, random objects with cores of lengths $14$ to $30$ for $L=8$, resp.\ $L=20$. Without the lemmas, the scattering of every ordered pair of sequences of a set $\mathcal U$ at both relative parities is tabulated, and the properties (C0)--(C3) are applied to the bare points until nothing is added, points outside $\mathcal U$ being discarded, so that the result is a subset of $C_L$. It is all of $\mathcal U$ when $\mathcal U$ consists of the $153$ sequences with cores of length at most $8$ and $1\le L\le6$ ($23{,}850$ simulations), of the $818$ sequences with cores of length at most $11$ and $L\in\{4,6,8,10,12\}$ ($682{,}120$ simulations), and of the $1434$ sequences with cores of length at most $12$ and $L\in\{4,6,8,10,12,13\}$ ($2{,}096{,}356$ simulations). When (C0) is left out, the result for the $267$ sequences with cores of length at most $9$ consists of the four bare points, in accordance with the paragraph after Definition~\ref{cl:def:CL}.
